% !PRINTED 482
\weberpage{482}
%<<CONT
so wird
\weq{21}{f_2(\omega) = \frac{e^{-\frac{r\pi i}{24}}\sqrt2}{x}, \quad f_2\left(\frac{\omega}{n}\right) = e^{-\frac{r\pi i}{24}}x,}
und diese Werte hat man für $u_1$, $v_1$ in das System II., § 73 zu
substituieren, um eine Gleichung für $x$ zu erhalten.
Indem man für $r$ die verschiedenen zulässigen Werte setzt,
bekommt man aus (16):
\[
\begin{aligned}
n &= 3, & m &= 6,\ 5,\ 2,\\
n &= 5, & m &= 10,\ 9,\ 6,\ 1,\\
n &= 7, & m &= 14,\ 13,\ 10,\ 5,\\
n &= 11, & m &= 22,\ 21,\ 18,\ 13,\ 6,\\
n &= 13, & m &= 26,\ 25,\ 22,\ 17,\ 10,\ 1,\\
n &= 17, & m &= 34,\ 33,\ 30,\ 25,\ 18,\ 9,\\
n &= 19, & m &= 38,\ 37,\ 34,\ 29,\ 22,\ 13,\ 2.
\end{aligned}
\]
Als einfaches Beispiel wählen wir $n = 7$ und erhalten aus
§ 73, II.:
\[ \frac{x^8}{4} + \frac{4}{x^8} = 7 + 4\sqrt2 \cos\frac{r\pi}{4}. \]
Daraus ergibt sich für $r = 0$:
\weq{22}{\frac{x^2}{\sqrt2} + \frac{\sqrt2}{x^2} = 1+\sqrt2, \quad x = f_1(\sqrt{-14}),}
für $r=1$ das bereits bekannte
\weq{23}{f(\sqrt{-13})^4 = 3+\sqrt{13},}
für $r=2$:
\weq{24}{f_1(\sqrt{-10})^2 = \frac{1+\sqrt5}{\sqrt2}.}
Als zweites Beispiel nehmen wir noch $n = 13$ und erhalten:
\weq{25}{A_1 = \frac{\sqrt2}{x^2} - \frac{x^2}{\sqrt2}, \quad B_1 = 16\cos\frac{r\pi}{2};}
für $r = 0, 4$, also $m = 26, 10$ ergibt sich hieraus $B_1 = 16$, und
folglich nach § 73:
\[ A_1^7 - 6A_1^5 + A_1^3 + 20A_1 + 16 = 0, \]
während für $r=2$, also $m=22$ in dieser Gleichung $A_1$ in
$-A_1$ zu verwandeln ist.
Man findet aber leicht die Zerlegung:
\weq{26}{A_1^7 - 6A_1^5 + A_1^3 + 20A_1 + 16 = (A_1+1)^2(A_1-2)^2(A_1^3+2A_1^2+A_1+4).}
